% 源分章 PDF 第 16—18 页；第 12.5 节及表 12.1—12.2。
\section{模拟}

下面的模拟研究旨在：(1) 进一步说明不同因果参数之间的差别；(2) 展示 TMLE 的实现；(3) 了解参数的指定如何影响估计量的精度及其所达到的功效。我们着重考虑随机试验，以说明设计阶段采用配对匹配，以及分析阶段进行调整，可能带来的效率提升。全部模拟均在 R 中进行（R Development Core Team，2016）。完整 R 代码见 Balzer 等（2016c）。

考虑单位 $i=\{1,\ldots,n\}$ 的如下数据生成过程。首先，从标准正态分布中抽取背景误差 $U_{Y,i}$。随后，从均值均为 0、标准差均为 1 的多元正态分布中生成五个基线协变量。前两个协变量 $(W1_i,W2_i)$ 的相关性为 0，后三个协变量 $(W3_i,W4_i,W5_i)$ 的相关性为 0.65。将暴露 $A_i$ 随机分配，使干预单位与对照单位数量相等。回顾一下，$A_i$ 是一个二元指示变量：单位被随机分配至干预组时为 1，被分配至对照组时为 0。在未匹配试验中，将 $n/2$ 个单位随机分配至干预组，其余单位分配至对照组。在匹配试验中，使用非二分匹配算法 \texttt{nbpMatch}（Beck 等，2016），根据 $\{W1,W4,W5\}$ 对单位进行配对。结局 $Y_i$ 按下式生成：
\[
\begin{aligned}
Y_i=\logit^{-1}\bigl[&A_i+0.75W1_i+0.75W2_i+1.25W3_i+U_{Y,i}\\
&+0.75A_iW1_i-0.5A_iW2_i-A_iU_{Y,i}\bigr]/5.
\end{aligned}
\]
还通过干预设置 $A_i=a$，生成反事实结局 $Y_{a,i}$。对于样本量 $n=\{30,50\}$，将这一数据生成过程重复 5000 次。SATE 的真值按每个样本中反事实结局之差的平均计算；PATE 的真值则通过在一个包含 500,000 个单位的总体中，对反事实结局之差取平均而计算。在该总体中，观测结局 $Y$ 与基线协变量之间的相关性从弱到中等：与 $W1$ 为 0.5，与 $W2$ 为 0.2，与 $W3$ 为 0.6，与 $W4$ 和 $W5$ 均为 0.4。

我们比较未调整估计量与采用两种结局回归初始估计方法的 TMLE 的表现。具体而言，先使用逻辑回归估计 $\bar Q_0(A,W)$，其中纳入暴露 $A$、协变量 $W1$ 以及交互项 $A*W1$。我们还使用超级学习器估计 $\bar Q_0(A,W)$，候选库包含所有如下形式的逻辑回归：纳入暴露 $A$、一个协变量以及二者的交互项。未调整估计量可视为 TMLE 的一个特例，其中 $\bar Q_n(A,W)=\EE_n(Y\mid A)$。推断基于估计影响曲线和 Student 的 $t$ 分布。我们构造 Wald 型 95\% 置信区间，并检验平均效应为零的原假设。

\noindent\textbf{结果。}表 12.1 汇总了 5000 次模拟试验中的参数取值。回顾一下，SATE 的真值依赖于纳入研究的单位，而总体仅有一个 PATE 真值。样本效应从 0.17\% 到 5.94\% 不等，均值为 2.97\%；总体效应恒为 2.98\%。如预期，SATE 的变异随样本量增大而降低。

\begin{table}[htbp]
\centering
\caption{样本量为 $n=\{30,50\}$ 时，5000 次模拟的因果参数汇总（单位：\%）}
\label{tab:ch12-parameters}
\begingroup
\small
\setlength{\tabcolsep}{4pt}
\begin{tabular}{@{}lrrrrrrrr@{}}
\toprule
& \multicolumn{4}{c}{SATE} & \multicolumn{4}{c}{PATE} \\
\cmidrule(lr){2-5}\cmidrule(l){6-9}
& 最小值 & 均值 & 最大值 & 方差 & 最小值 & 均值 & 最大值 & 方差 \\
\midrule
$n=30$ & 0.17 & 2.97 & 5.94 & $6.5\mathrm E{-3}$ & 2.98 & 2.98 & 2.98 & 0 \\
$n=50$ & 0.18 & 2.96 & 5.14 & $4.2\mathrm E{-3}$ & 2.98 & 2.98 & 2.98 & 0 \\
\bottomrule
\end{tabular}
\endgroup
\end{table}

表 12.2 展示了各估计量的表现。具体而言，我们将偏差定义为点估计与（样本特定）真值之差的平均，将标准差定义为估计量相对于其目标的方差的平方根，并给出均方误差（mean squared error，MSE）。还给出相对均方误差（relative MSE，rMSE），即某个估计量的 MSE 除以未匹配试验中总体效应未调整估计量的 MSE。此外，表中还列出所达到的功效，即原假设为假时被拒绝的比例，以及 95\% 置信区间覆盖率。

\begin{table}[htbp]
\centering
\caption{5000 次模拟中估计量表现的汇总}
\label{tab:ch12-performance}
\begingroup
\small
\setlength{\tabcolsep}{3pt}
\begin{adjustbox}{max width=\textwidth}
\begin{tabular}{@{}llrrrrrr@{}}
\toprule
目标与设计 & 估计量 & 偏差 & 标准差 & MSE & rMSE & 功效 & 覆盖率 \\
\midrule
\multicolumn{8}{c}{样本量 $n=30$} \\
PATE，未匹配 & 未调整 & $2.3\mathrm E{-4}$ & $2.2\mathrm E{-2}$ & $4.8\mathrm E{-4}$ & 1.00 & 0.27 & 0.95 \\
 & TMLE & $6.8\mathrm E{-4}$ & $1.9\mathrm E{-2}$ & $3.6\mathrm E{-4}$ & 0.75 & 0.36 & 0.94 \\
 & TMLE+SL & $2.9\mathrm E{-4}$ & $1.6\mathrm E{-2}$ & $2.6\mathrm E{-4}$ & 0.55 & 0.48 & 0.93 \\
SATE，未匹配 & 未调整 & $3.1\mathrm E{-4}$ & $2.0\mathrm E{-2}$ & $4.2\mathrm E{-4}$ & 0.88 & 0.27 & 0.96 \\
 & TMLE & $7.5\mathrm E{-4}$ & $1.7\mathrm E{-2}$ & $3.0\mathrm E{-4}$ & 0.63 & 0.39 & 0.95 \\
 & TMLE+SL & $3.7\mathrm E{-4}$ & $1.4\mathrm E{-2}$ & $2.0\mathrm E{-4}$ & 0.42 & 0.52 & 0.95 \\
SATE，已匹配 & 未调整 & $5.4\mathrm E{-5}$ & $1.5\mathrm E{-2}$ & $2.2\mathrm E{-4}$ & 0.46 & 0.37 & 0.98 \\
 & TMLE & $3.7\mathrm E{-4}$ & $1.4\mathrm E{-2}$ & $2.1\mathrm E{-4}$ & 0.43 & 0.44 & 0.97 \\
 & TMLE+SL & $1.3\mathrm E{-4}$ & $1.1\mathrm E{-2}$ & $1.3\mathrm E{-4}$ & 0.27 & 0.58 & 0.97 \\
\midrule
\multicolumn{8}{c}{样本量 $n=50$} \\
PATE，未匹配 & 未调整 & $-1.3\mathrm E{-4}$ & $1.7\mathrm E{-2}$ & $3.0\mathrm E{-4}$ & 1.00 & 0.41 & 0.94 \\
 & TMLE & $1.1\mathrm E{-4}$ & $1.5\mathrm E{-2}$ & $2.2\mathrm E{-4}$ & 0.75 & 0.53 & 0.94 \\
 & TMLE+SL & $-3.1\mathrm E{-6}$ & $1.2\mathrm E{-2}$ & $1.6\mathrm E{-4}$ & 0.53 & 0.68 & 0.94 \\
SATE，未匹配 & 未调整 & $4.8\mathrm E{-5}$ & $1.6\mathrm E{-2}$ & $2.5\mathrm E{-4}$ & 0.86 & 0.41 & 0.96 \\
 & TMLE & $2.9\mathrm E{-4}$ & $1.3\mathrm E{-2}$ & $1.8\mathrm E{-4}$ & 0.60 & 0.55 & 0.96 \\
 & TMLE+SL & $1.8\mathrm E{-4}$ & $1.1\mathrm E{-2}$ & $1.1\mathrm E{-4}$ & 0.38 & 0.70 & 0.97 \\
SATE，已匹配 & 未调整 & $-1.8\mathrm E{-4}$ & $1.1\mathrm E{-2}$ & $1.1\mathrm E{-4}$ & 0.38 & 0.59 & 0.98 \\
 & TMLE & $-1.6\mathrm E{-4}$ & $1.0\mathrm E{-2}$ & $1.1\mathrm E{-4}$ & 0.36 & 0.66 & 0.97 \\
 & TMLE+SL & $-5.7\mathrm E{-5}$ & $8.2\mathrm E{-3}$ & $6.7\mathrm E{-5}$ & 0.23 & 0.81 & 0.98 \\
\bottomrule
\end{tabular}
\end{adjustbox}
\par\smallskip
\begin{minipage}{\textwidth}
各行给出目标参数、研究设计和估计量：未调整估计量、采用逻辑回归的 TMLE，以及采用超级学习器的 TMLE（“TMLE+SL”）。
\end{minipage}
\endgroup
\end{table}

如预期，所有估计量均无偏。在随机试验中，$\bar Q_0(A,W)$ 回归模型的误设不会带来偏差风险（例如 Tsiatis 等，2008；Moore 和 van der Laan，2009b；Rosenblum 和 van der Laan，2010b）。同样如预期，估计量的精度随样本量增大和进行调整而提高（例如 Fisher，1932；Cochran，1957；Cox 和 McCullagh，1982；Tsiatis 等，2008；Moore 和 van der Laan，2009b；Rosenblum 和 van der Laan，2010b）。例如，考虑在一个有 $n=30$ 个单位且未匹配的试验中估计总体效应。未调整估计量的标准误为 $2.2*10^{-2}$；调整一个协变量后，标准误为 $1.9*10^{-2}$。通过超级学习器纳入条件均值 $\bar Q_0(A,W)$ 的数据自适应估计后，标准误进一步降至 $1.6*10^{-2}$。配对匹配也如预期提高了精度（Imai 等，2009；van der Laan 等，2013a；Balzer 等，2015）。对于 SATE，在 $n=30$ 和 $n=50$ 时，未匹配试验中未调整估计量的标准误，分别是配对匹配试验中相应估计量的 1.38 倍和 1.49 倍。

对于所有估计算法和样本量，目标参数的指定对精度和功效都有显著影响。正如理论预测，PATE 的未调整估计量具有最高方差。当单位数为 $n=50$ 时，该估计量用于 PATE 的 MSE，是未匹配试验中结合超级学习器的 SATE TMLE 的 2.62 倍，是匹配试验中结合超级学习器的 SATE TMLE 的 4.42 倍。在有限样本模拟中，渐近保守的方差估计量对样本效应推断的影响十分明显。多数情形下，SATE 估计量的标准差被高估，置信区间覆盖率大于或等于 95\% 的名义水平。尽管方差估计量保守，样本效应 TMLE 的功效仍高于相同的总体效应 TMLE。当 $n=30$ 时，结合超级学习器的 TMLE 对总体效应达到的功效为 48\%，对未匹配样本效应为 52\%，对配对匹配后的样本效应为 58\%。当 $n=50$ 时，相应功效分别为 68\%、70\% 和 81\%。值得注意的是，在未匹配试验中，两个参数的未调整估计量具有相同功效。这是因为所估计影响曲线的 $D_W(O)$ 分量及其方差为零：
\[
\EE_n(Y\mid A=1)-\EE_n(Y\mid A=0)-\Psi_{n,\mathrm{unadj}}(P_n)=0,
\]
其中，$\EE_n(Y\mid A)$ 表示特定处理下的均值。因此，使用未调整估计量，会失去将 SATE 指定为推断目标原本可能带来的全部功效提升。相比之下，使用超级学习器的 TMLE 能够更好地拟合结局回归 $\bar Q_0(A,W)$，并获得不那么保守的方差估计量，因此达到最高功效。需要指出，在小型试验中，例如 $n\leq30$ 的早期临床试验或整群随机试验，精确估计 $\bar Q_0(A,W)$ 很可能具有挑战。实际中，许多基线协变量能够预测结局，但调整过多协变量可能导致过拟合。第 13 章将讨论一种利用交叉验证、从预先指定候选库中以数据自适应方式选择最优调整集的程序。
\FloatBarrier
